\section{Logarithm}

In this section, we will discuss formal logarithm on formal group law. Consider
integrating an invariant differential form, we expect that it will give us a homomorphism
to the additive group. However, integration will introduce denominators. In the case
of characteristic $0$ or torsion free, things work as we expect.

From now on, we fix a torsion free ring $R$, let $K = R \otimes \QQ$.

\begin{definition}
Let $F$ be a formal group law over $R$, $\omega(X)$ be its normalized invariant differential.
We denote $\omega(X) = (1 + a_1 X + a_2 X ^ 2 + \cdots) dX$, then the formal logarithm of $F$
is defined to be
    \[ \log_F (X) := \int \omega = X + \frac{a_1}{2} X^ 2 + \frac{a_2}{3} X ^ 3 + \cdots \in K [\![X]\!] \]
Since the coefficient of $X$ in $\log_F(X)$ is $1$, which is invertible in $K$, then
$\exp_F (X)$ is defined to be the unique power series over $K$ satisfying
    \[ \exp_F( \log_F (X)) = X. \]
\end{definition}

\begin{example}
    The normalized invariant differential on $\bfG_a$ is $dX$, then we have
    \[ \log_{\bfG_a} (X) = X, \ \exp_{\bfG_a}(X) = X. \]
    The normalized invariant differential on $\bfG_m$ is $(1 + X) ^ {-1} dX$, then we have
    \[ \log_{\bfG_m} (X) = \sum_{n=1}^{\infty} \frac{(-1)^{n-1}X ^ n}{n}  \]
    \[ \exp_{\bfG_m} (X) = \sum_{n=1}^{\infty} \frac{X^n}{n!} \]
\end{example}

As we expect, the formal logarithm on $F$ is a homomorphism from $F$ to $\bfG_a$.
However, since $\log_F$ introduce denominators, this is no longer a homomorphism over $R$,
it is a homomorphism over $K$. Moreover, it is a isomorphism over $K$.

\begin{proposition}
    Let $F$ be a formal group law over $R$, then $\log_F(X)$ is a formal group isomorphism from
    $F$ to $\bfG_a$ over $K$.
\end{proposition}

\begin{proof}
    Let $\omega(X)$ be the normalized invariant differential for $F$. Then we have
    \[\omega \circ F = \omega .\]
    Integrating this, we have
    \[\log_F(F(X,Y)) = \log_F(X) + C(Y). \]
    where $C(Y) \in K[\![Y]\!]$. Let $X = 0$ in above, we have
    \[ C(Y) = \log_F (0) + C (Y) = \log_F(F(0,Y)) = \log_F (Y). \]
    Then we have that
    \[\log_F(F(X,Y)) = \log_F(X) + \log_F(Y). \]
    which proves that $\log_F(X)$ is a homomorphism from $F$ to $\bfG_a$. Furthermore,
    since $\exp_F$ is its inverse, then $\exp_F$ is homomorphism from $\bfG_a$ to $F$.
    Therefore, $\log_F$ is an isomorphism.
\end{proof}

\begin{corollary}
    Let $R$ be a torsion free ring, then every formal group law over $R$ is commutative.
\end{corollary}

\begin{proof}
    The assumption that $R$ is torsion free guarantee the existence of $\log_F$ and $\exp_F$.
    Then
        \[F(X,Y) = \exp_F(\log_F(X) + \log_F(Y)). \]
    Therefore, we have that $F(X,Y) = F(Y,X)$.
\end{proof}

There is some more discussion of the form of coefficient of $\log_F$ and $\exp_F$ in
\cite{silverman2009arithmetic}.